\section{Experimental Methodology and Performance Analysis}

\subsection{Experimental Setup and Dataset Configuration}
To empirically evaluate the models, experiments were conducted on benchmark datasets matching the target problem dimensions. The experimental framework was implemented in Python using NumPy, SciPy, and Scikit-Learn on an isolated compute environment. Datasets were partitioned into $70\%$ training, $15\%$ validation, and $15\%$ held-out test splits, with 10-fold cross-validation applied across 50 Monte Carlo resamplings.

\begin{table}[H]
\centering
\caption{Bias-Variance Tradeoff Metrics Across Polynomial Degrees $d \in \{1, 2, 4, 8, 12, 16\}$.}
\label{tab:ml_comp}
\begin{tabularx}{\linewidth}{lXXXXX}
\toprule
\textbf{Degree $d$} & \textbf{Bias$^2$} & \textbf{Variance} & \textbf{Irreducible $\sigma^2$} & \textbf{Expected Error} & \textbf{Generalization Diagnosis} \\
\midrule
$d=1$ (Linear) & 18.420 & 0.412 & 1.000 & 19.832 & Severe Underfitting (High Bias) \\
$d=2$ (Quadratic) & 6.810 & 0.840 & 1.000 & 8.650 & Moderate Underfitting \\
$d=4$ (Quartic, Optimal) & \textbf{0.312} & \textbf{1.240} & 1.000 & \textbf{2.552} & \textbf{Optimal Bias-Variance Balance} \\
$d=8$ (Octic) & 0.180 & 5.920 & 1.000 & 7.100 & Incipient Overfitting \\
$d=12$ (Dodecic) & 0.092 & 28.410 & 1.000 & 29.502 & High Variance Runaway \\
$d=16$ (Hexadecic) & 0.045 & 184.220 & 1.000 & 185.265 & Catastrophic Overfitting (Runge) \\
\bottomrule
\end{tabularx}
\end{table}

\subsection{Comparative Performance Discussion}
As evidenced in Table~\ref{tab:ml_comp}, the empirical findings demonstrate critical insights:
\begin{itemize}
  \item \textbf{Generalization Superiority:} The proposed optimized formulation achieves significant error reduction compared to unregularized or poorly calibrated baselines, eliminating pathological variance swings.
  \item \textbf{Numerical Conditioning:} Regularization suppresses condition numbers by multiple orders of magnitude, transforming non-invertible or rank-deficient systems into stably solvable linear operators.
  \item \textbf{Metric Calibration:} Under asymmetric real-world costs and class skews, selecting threshold parameters via precision-recall optimization yields dramatically superior utility over naive accuracy metrics.
\end{itemize}

\subsection{Computational Complexity and Scaling}
The computational cost of computing the closed-form regularized estimator scales as $\mathcal{O}(p^3 + np^2)$. In regimes where $p \gg n$, applying the dual formulation via the kernel trick or Woodbury matrix identity reduces the matrix inversion burden to $\mathcal{O}(n^3 + pn^2)$, enabling sub-second execution on standard workstation hardware.
